% =====================================================================
\section{Limits and outlook}
\label{sec:limiti}
% =====================================================================

\begin{itemize}
  \item \textbf{Adoption.} 015 is used in real networks, 018 almost only in
        testbeds and prototypes~\cite{madqci2024}.
  \item \textbf{Incomplete multi-domain support.} The normative scenario is
        QKD$+$OTN; the QKD$+$QKD case requires extensions like those of Sim
        et al.\ and coordination with GS~QKD~020 for key hand-over between
        KMSs.
  \item \textbf{Metrics.} ESKR is an instantaneous indicator and accounts
        neither for key already accumulated in buffers nor for its
        dynamics. Laschet et al.~\cite{laschet2026} propose richer
        indicators for path selection.
  \item \textbf{Topology confidentiality.} Exposing the internal topology
        to the orchestrator is acceptable within one operator, less so
        between different operators; inter-domain routing schemes that
        preserve confidentiality are needed~\cite{laschet2026,interconnect2026}.
  \item \textbf{Interface security.} The standard does not specify how the
        orchestrator–controller channel is protected, leaving it to
        implementations.
  \item \textbf{Algorithms.} Path computation is not standardised: this
        helps innovation~\cite{routing2024,mendez2026} but makes behaviour
        less predictable in multi-vendor scenarios.
\end{itemize}

% =====================================================================
\section{Conclusions}
\label{sec:conclusioni}
% =====================================================================

Read in sequence, ETSI's work on QKD control follows a clear logic.
GS~QKD~002 showed that useful QKD is a network, with key management,
differentiated service requirements and integration into the operator's
optical network. GS~QKD~015 answered with the SD-QKD node abstraction:
modules as interfaces, a single key manager, physical and virtual links,
applications with QoS, and a controller that registers, plans and
configures, all described in YANG. GS~QKD~018 added the orchestration
layer: each domain becomes an abstract graph, key resources are monitored
through the inventory and provisioning follows a scheme in which the
controller proposes and the orchestrator decides. The design mirrors ACTN,
reuses YANG, NETCONF and RESTCONF and keeps the essential security
principle: keys never leave the nodes. The KOREN, Politecnico di Milano and
MadQCI experiences show that the scheme works but that QKD orchestration is
still young. The next step is to truly extend it to multi-operator
settings, with better metrics and confidentiality guarantees between
domains.
